\documentclass[11pt]{article}
\usepackage{manual_docs_style}
\externaldocument{build/environment_profiles}[environment_profiles.pdf]
\externaldocument{build/candidate_generation}[candidate_generation.pdf]
\externaldocument{build/simulation_stage}[simulation_stage.pdf]
\externaldocument{overleaf_paper/build/main}[overleaf_paper/archive/compiled/main.pdf]

\begin{document}

\docTitle{Stage: Cheap Filter Metrics}
\phantomsection\label{docstart:cheap_filter_metrics_stage}

\subsection*{Overview}
This stage computes low-cost detectability filters for materialized intervention runs.
For each intervention, it compares the trajectory against the matched \hyperref[sec:candidate-generation-baseline-recipe]{baseline} and the matched \hyperref[sec:candidate-generation-time0-baseline-recipe]{time0 baseline}.
The generic metric is a per-signal RMS difference, normalized over the valid aligned samples for that signal.
The stage can also call a \termSimulator-specific cheap detectability hook to reject cases that are known to be indiguishable between classes.

Thresholds are channel-specific.
They determine whether valid trajectories are distinguishable at or after the intervention.
Pre-intervention divergence is not a detectability-threshold concern: Stage Simulate first requires each intervention to match its family baseline within absolute tolerance \code{1e-3} at timestamps strictly before the intervention.
A mismatch there is recorded as a simulation-stage integrity failure, so the cheap-filter stage must not conceal or reinterpret it as detectability.
\agentProposedEdit{tracebench-docs-solvability-section-reference}{The hyperlink target should use the renamed label for the Solvability filtering subsection.}{Additional information can be found in the paper section on \hyperref[sec:detectability]{detectability and distinguishability}.}{Additional information can be found in the paper section on \hyperref[sec:solvability-filtering]{detectability and distinguishability}.}

\section*{Invocation}

The script is \code{workflows/metrics/compute_metrics.py}.

\begin{DocCode}
python workflows/metrics/compute_metrics.py \
  --model_dir SIMULATOR_DIRECTORY \
  --policy POLICY_ID \
  --root_dir ROOT_DIR \
  [--uuid UUID] \
  [--thresholds THR1 ...] \
  [--thresholds-file THRESHOLDS_JSON] \
  [--filter-profile legacy|strict_clean_v1] \
\end{DocCode}

Here \code{ROOT_DIR} is the run-artifact root produced by \hyperref[docstart:simulation_stage]{Stage: Simulate}.
It is the directory containing \code{model_record.json} and the per-run artifact directories.
During execution, the script displays a progress bar over the intervention runs in the selected policy.
\code{--thresholds} and \code{--thresholds-file} are mutually exclusive.
The threshold file may be either a direct signal-to-value JSON object or a provenance object containing a non-empty \code{thresholds} object.
The default profile is \code{legacy}.

\section*{Inputs}

For each \termSimulator, the stage reads:

\begin{itemize}
\item \code{plans/<policy_id>/}\hyperref[sec:candidate-generation-run-nodes]{\code{run_nodes.jsonl}} and \code{plans/<policy_id>/}\hyperref[sec:candidate-generation-run-edges]{\code{run_edges.jsonl}}.
\item  \code{<model_dir>/}:
\begin{DocCode}
<model_dir>/
    experiment_config.json
    plans/<policy_id>/run_nodes.jsonl
    plans/<policy_id>/run_edges.jsonl
    detectability_specific_environment.py # optional
    noise_adder.py
\end{DocCode}
\item \code{<run_root>/}, the run-artifact root:
\begin{DocCode}
<run_root>/
  model_record.json
  <run_id>/
    recipe.json
    data.parquet
    simulink_model.mdl
    features.json        # optional
    debug/               # optional
  <another_run_id>/
    ...
\end{DocCode}
\item \code{experiment_config.json}: used to read default thresholds when thresholds are not passed explicitly.
\item \code{detectability_specific_environment.py}: optional \termSimulator-specific detectability hook.
\end{itemize}

By default, \code{<run_root>} is:

\begin{DocCode}
models/simulink/<simulator_id>/runs/
\end{DocCode}

When Stage: Simulate is run with \code{--runs-root RUNS_ROOT}, \code{<run_root>} is:

\begin{DocCode}
<RUNS_ROOT>/<simulator_id>/runs/
\end{DocCode}

\section*{Outputs}

Each cheap-filter execution writes one metric-run directory:

\begin{DocCode}
<model_dir>/metrics/<policy_id>/<metric_run_id>/
  metric_run.json
  cheap_filter_metrics.jsonl
\end{DocCode}

\code{<model_dir>} is the path to a folder such \code{models/simulink/BallDrop}.
\code{metric_run_id} identifies the exact cheap-filter configuration that produced the metrics.
Downstream stages should consume the metric-run directory, not infer configuration from a UUID embedded in a filename.

\section*{Schema}

\code{metric_run.json} records the identity and configuration for one cheap-filter execution:

\begin{DocCode}
{
  "record_type": "metric_run",
  "metric_run_id": "cheap_filter_8f4c9d2e",
  "metric_type": "cheap_filter",
  "created_at": "2026-06-23T12:00:00Z",
  "script": "workflows/metrics/compute_metrics.py",
  "model": "BallDrop",
  "policy_id": "production_v1",
  "model_dir": "models/simulink/BallDrop",
  "run_root": "models/simulink/BallDrop/runs",
  "inputs": {
    "run_nodes": "plans/production_v1/run_nodes.jsonl",
    "run_edges": "plans/production_v1/run_edges.jsonl",
    "model_record": "runs/model_record.json"
  },
  "thresholds": {
    "height": 0.02,
    "velocity": 0.05
  },
  "detectability_specific_environment": "detectability_specific_environment.py"
}
\end{DocCode}

\code{cheap_filter_metrics.jsonl} is JSON Lines.
Each intervention writes one run record:

\begin{DocCode}
{
  "record_type": "run",
  "metric_run_id": "cheap_filter_8f4c9d2e",
  "policy_id": "production_v1",
  "family_id": "fam_...",
  "run_id": "run_intervention_...",
  "kind": "intervention",
  "ground_truth_label": "gravity",
  "baseline_run_id": "run_baseline_...",
  "time0_run_id": "run_time0_...",
  "changed_parameter": "gravity",
  "filter_profile": "strict_clean_v1",
  "filter_pass": true,
  "cheap_filter_pass": true,
  "strict_clean_pass": true,
  "linked_run_status": {
    "baseline": "success",
    "intervention": "success",
    "time0_baseline": "success"
  },
  "baseline_rms_difference": {
    "height": 0.013,
    "velocity": 0.041
  },
  "time0_rms_difference": {
    "height": 0.020,
    "velocity": 0.058
  },
  "detectability": {
    "environment_specific_detectability": true,
    "vs_baseline": {
      "generic_numeric_detectability": true,
      "detectability_output": {
        "first_diff": [1.23],
        "mean_SNR": [31.4]
      }
    },
    "vs_time0_baseline": {
      "generic_numeric_detectability": true,
      "detectability_output": {
        "first_diff": [1.20],
        "mean_SNR": [29.8]
      }
    }
  },
  "failure_reasons": []
}
\end{DocCode}


\section*{Enviroment specific Detectability}
This uses the \termSimulator-specific function \code{is_detectable} as defined here:

\section*{Pass Rules}

\code{filter_pass} is the result selected by \code{filter_profile}.
Both profile-specific booleans are always written, so changing the selected profile does not erase the legacy result.

An intervention has \code{cheap_filter_pass = true} only if:

\begin{itemize}
\item it has status \code{success} in \code{model_record.json};
\item it is numerically valid and has all required observable signals;
\item \code{generic_numeric_detectability} is true for both \code{vs_baseline} and \code{vs_time0_baseline}.
\end{itemize}
Note for generating valid deep learning samples we only care about \code{generic_numeric_detectability} is true for both \code{vs_baseline} and \code{vs_time0_baseline}.

An intervention has \code{strict_clean_pass = true} only if all three linked runs have status \code{success}; all three trajectories are present and numerically valid; generic numeric detectability passes against both the family baseline and time-zero baseline; the simulator-specific hook passes against the family baseline; and each comparison reports a finite first-difference time at or after the declared intervention time.

Strict-clean thresholds can be calibrated with \code{workflows/metrics/calibrate_strict_clean_thresholds.py}.
The campaign calibration applies the simulator's \code{high} noise profile with seeds 0--4 to successful pilot trajectories, computes each signal's full-trajectory L2 distance from its clean trajectory, and writes \code{1.10} times the empirical 99th percentile together with input hashes and sample counts.

\section*{Validation Rules}

\begin{itemize}
\item Every metric record must refer to a \code{run_id} present in \code{run_nodes.jsonl}.
\item Every intervention metric must have a matched baseline and time-zero/static-change run from \code{run_edges.jsonl}.
\item Only successful materialized runs may have \code{cheap_filter_pass = true}.
\item The cheap-filter stage must not change run recipes or recipe hashes.
\item Re-running this stage with the same \code{metric_run_id} may overwrite the existing metric-run directory.
\end{itemize}

\end{document}
